	\section{Conclusiones}
	
	El presente Trabajo Terminal permitió desarrollar e implementar un prototipo funcional de reconocimiento de señas de la LSM en dispositivos móviles Android, integrando técnicas de visión por computadora, aprendizaje profundo e inferencia local optimizada para entornos móviles. A lo largo de aproximadamente un año de trabajo se llevaron a cabo múltiples etapas de investigación, experimentación, diseño, desarrollo, pruebas y validación, permitiendo no solamente construir un prototipo funcional, sino también comprender de manera profunda las limitaciones y retos reales asociados al reconocimiento automático de señas en dispositivos móviles de propósito general.
	
	La primer parte del proyecto se exploraron distintas arquitecturas de CNN basadas en procesamiento de imágenes RGB completas, incluyendo técnicas de segmentación y escalas de grises. Conforme avanzaron las pruebas experimentales, se observó que dichos enfoques presentaban una fuerte dependencia respecto a condiciones de iluminación, fondos complejos, posiciones de la mano y variaciones entre usuarios, dificultando considerablemente la capacidad de generalización del sistema en escenarios reales.
	
	Como resultado de estas observaciones, el proyecto evolucionó hacia un enfoque basado en extracción de características geométricas mediante MediaPipe Hands, utilizando landmarks de la mano como representación principal de las señas. Esta decisión representó uno de los cambios más importantes dentro del desarrollo del prototipo, ya que permitió reducir significativamente la complejidad del problema al trabajar únicamente con coordenadas normalizadas de puntos clave de la mano en lugar de imágenes completas. Gracias a ello, fue posible disminuir el costo computacional, mejorar la velocidad de inferencia y obtener un sistema considerablemente más robusto para ejecución en dispositivos móviles.
	
	A lo largo de las diferentes etapas de entrenamiento y pruebas se realizaron múltiples experimentos relacionados con normalización de coordenadas, variaciones en la orientación del dispositivo, tamaños de entrada, técnicas de suavizado temporal y estrategias de confirmación de predicciones. Estas pruebas permitieron identificar comportamientos importantes que difícilmente podrían haberse anticipado únicamente desde el enfoque teórico. Uno de los hallazgos más relevantes fue la diferencia existente entre las inferencias realizadas en computadora y aquellas ejecutadas directamente en Android. Aunque el modelo mostraba resultados altamente estables en entorno de escritorio, durante las pruebas móviles comenzaron a observarse inconsistencias relacionadas con el pipeline interno de procesamiento de imagen, la orientación del dispositivo y el comportamiento específico de MediaPipe dentro de Android.
	
	Particularmente se observó que la orientación horizontal del dispositivo proporcionaba resultados considerablemente más estables que la orientación vertical, reduciendo errores de clasificación y mejorando la consistencia de las predicciones. Asimismo, se comprobó que factores como la iluminación, la velocidad del movimiento, el encuadre de la mano y entre otros factores tienen un impacto directo sobre la precisión del sistema. Estas observaciones llevaron a implementar mecanismos adicionales como regiones de interés, suavizado temporal y máquinas de estados para disminuir cambios o variaciones y reducir predicciones erráticas entre frames consecutivos.
	
	Otro aspecto importante identificado en el desarrollo fue que, aunque MediaPipe Hands facilita enormemente la extracción de landmarks, existen diferencias sutiles entre la implementación utilizada en computadora y aquella integrada en dispositivos Android. Estas diferencias pueden generar pequeñas variaciones en las coordenadas obtenidas en la inferencia móvil, afectando la capacidad de generalización del modelo entrenado originalmente con datos capturados en entorno de escritorio. Dada esta situación una de las principales recomendaciones derivadas de este trabajo consiste en reconstruir completamente el dataset directamente desde dispositivos móviles Android, utilizando exactamente el mismo pipeline de captura, orientación, resolución y extracción de landmarks empleado en la inferencia final de la aplicación. Esto permitiría reducir inconsistencias entre entrenamiento e inferencia y potencialmente mejorar de forma considerable la estabilidad del prototipo en escenarios complejos.
	
	Se concluye que la calidad y diversidad del dataset representan uno de los factores más críticos dentro del reconocimiento automático de señas. Aunque el sistema logró resultados funcionales y prometedores, aún existen áreas importantes de mejora relacionadas con la incorporación de un mayor número de usuarios y diversidad entre ellos, condiciones de iluminación variables, fondos complejos y diferentes configuraciones de mano. Incrementar la diversidad de datos permitiría fortalecer la capacidad de generalización del modelo y acercarlo a condiciones reales de uso cotidiano.
	
	Hablando de arquitectura y rendimiento, el uso de TensorFlow Lite demostró ser una solución para inferencia móvil debido a su bajo consumo computacional y compatibilidad con dispositivos Android. El empleo de landmarks como representación principal permitió construir un prototipo ligero y eficiente, capaz de ejecutarse localmente sin depender de servidores externos o conexión a internet, aspecto particularmente importante desde el punto de vista de privacidad, accesibilidad y portabilidad.
	
	Más allá de los resultados técnicos obtenidos, este proyecto también permitió comprender la importancia social y humana asociada al desarrollo de herramientas tecnológicas orientadas a inclusión y accesibilidad. La LSM representa un medio de comunicación fundamental para la comunidad sorda y aunque la aplicación desarrollada aún corresponde a un prototipo académico, los resultados obtenidos demuestran el potencial que poseen las tecnologías visión por computadora para servir como base en futuras herramientas de apoyo a la comunicación.
	
	Este Trabajo Terminal establece una base sólida para futuras líneas de investigación y desarrollo. Entre los trabajos futuros más importantes destacan la incorporación de señas dinámicas mediante arquitecturas secuenciales como LSTM o transformers, el reconocimiento de palabras y frases completas, la expansión del diccionario de señas y la optimización multiplataforma hacia otros entornos móviles. En conjunto, los resultados obtenidos demuestran que el reconocimiento señas de la LSM en dispositivos móviles es un problema complejo y que aún existe un amplio potencial de mejora y expansión para futuras versiones del prototipo.
	
	\newpage
	
